\Note 1:6 {walking}
{\it AE 97} Walking in the Word refers to {\bf living}. {\bf Ways} refers to truths by which one 'walks' or lives.

\Note 1:11 {the Lord}
{\it AC 2921:5} Throughout the New Testament, God is referred to as {\bf the Lord} instead of Yahweh/Jehovah because a) if he had taken this name on people would not have believed him and b) he was not Yahweh/Jehovah until the last temptation (the crucifixion) because at that point he unified his Divine Essence with his Human Essence.

\Note 1:11 {right side}
{\it AE 298:13} The {\bf right side} refers to, in the context of men and angels in the Word, the wisdom and intelligence that they have from Divine good through Divine truth proceeding from the Lord. The angel therefore appeared to Zacharias in the wisdom and intelligence that they had from the Lord.

\Note 1:13 {Fear not}
{\it AE 80, AE 677:8} In the Word, 'the fearful' are those who are in no faith, as Zacharias is here. {\bf Fear} signifies in the Word various disturbances of mind arising from the influx of Divine Good and Truth. In the world, all those who come suddenly from their own life into a life which is in some degree spiritual, are at first afraid, but are renewed by the Lord. This renewal, or resuscitation, is effected by the Divine presence, accommodated to the person's own reception of good by means of truth, and is signified by {\bf Fear not}. It then leads to a new life of increased humility and adoration. {\bf i.e. Zacharias (those without faith) is renewed by the Lord by means of whatever good and truth he possessed.}

\Note 1:14 {joy and gladness}
{\it AE 660:4, AR 507} {\bf Joy and Gladness} in the Word signify the delights of the will and the understanding. {\bf Joy} signifies pleasure from good and {\bf Gladness} signifies pleasure from truth. These two phrases go together throughout the Word and indicate the heavenly marriage of good and truth. {\bf Zacharias (the fearful; without faith) will be or has been renewed by the Lord, and will go on to experience these delights.}

\Note 1:15 {Holy Ghost}
{\it AE 710:31} That John the Baptist was to be filled with the {\bf Holy Ghost} signifies that he was to represent the Lord in relation to the Word, as Elijah did. In the Word, which is Divine truth, there is everywhere the marriage of Divine good and Divine truth, and Divine good united with Divine truth is the Divine proceeding from the Lord, which is called the Holy Spirit.

\Note 1:17 {the spirit and power of Elijah}
{\it AC 6752} As in other places in the Word, Elijah represents the Lord in relation to the Word. Elijah represents the Word as to the Prophetic books, Moses represents the Word as to the Historical books/the Law. For John the Baptist to come 'in the spirit and power of Elijah' shows that he was to take up this mantle as well. 'Spirit' corresponds to Divine Truth, while 'Power' corresponds to Divine Good. 

\Note 1:17 {the hearts of the fathers to the children}
{\it AC 3703:9} {\bf The hearts of the fathers and children} signify the the goods and truths of the church, which the Lord was about to restore.

\Note 1:19 {that stand in the presence of God}
{\it AE 414, AR 366} To stand in the presence of God signifies one's ability to receive Divine Truth and live by it. The evil are not able to stand before God in heaven because they find the influx of Divine Good and Truth unbearable, but the good are able to withstand that influx without problem.

\Note 1:25 {he looked on me, to take away my reproach among men}
{\it AE 721:16} Since a "mother" signifies the church, and "sons and daughters" its truths and goods, and in the ancient churches, and afterwards in the Jewish church, all things were representative or significative, it was a reproach and disgrace for women to be barren, as in Gen. 30 where Rachel's reproach is taken away when she bears Joseph.

\Note 1:32 {Son of the Highest}
{\it Lord 29:2} The fact that the Lord is called the {\bf Sone of the Highest} is evidence of the fact that he was the Divine as to his spirit, while while his Human was from Mary. Where one part is Divine, the other must be as well, showing that he united his Divine to his Human. This also shows that the Father was in the Son, and vice versa, i.e. there is no separation of entities in God. The Divine Soul and Human physical cannot be separated any more than any mortal man's soul and body.

\Note 1:33 {reign over the house of Jacob}
{\it AC 3305:3} By {\bf Jacob} in the word is signified the doctrine of natural truth. Here {\bf the house of Jacob} was not used to mean the Jewish nation or people, because the Lord's kingdom included not merely that people, but people throughout the world who have faith in Him, and from faith have charity.

\Note 1:33 {of his kingdom there shall be no end}
{\it AE 1104:5} While churches come and go, as has happened four times on Earth already, when the Word refers to things of the Lord or Heaven in terms of ages or epochs, it is referring to what is timeless and eternal in those things, namely Good and Truth. The Lord's kingdom (his Good and Truth) actually are without end; it is only humans and their receptivity that limit this.

\Note 1:35 {The Holy Ghost shall come upon thee, and the power of the Highest shall overshadow thee}
{\it Ath 45-46, Can 17} The Lord's spirit was from the Father, not Mary or Joseph i.e. he was fully God. {\bf The Holy Ghost} signifies Divine Truth, while the {\bf power of the Highest} signifies Divine Good. In the Lord, these two were initially separate, as in man, but were united in time, just as they are united in man through regeneration.

\Note 1:35 {Son of God}
{\it AE 1069:2, TCR 342} God's Essence is his Divine Love, from which proceeds his Divine Good and Truth (his Existence, or manifestation). Since the Lord is the Word in multiple ways and because he proceeds from the Divine Essence, he is called the {\bf Son of God}. Since there is no way to have faith in Christ without acknowledging his Divinity, this is a core principle of Christianity. Without it, we would need to have faith in a man, which is a damaging doctrine, seen paralleled by the crown of thorns, and the taunting that Christ received on the cross to save himself.

\Note 1:41 {the babe leaped in her womb}
{\it AE 710:31} John leaping in the womb at the greeting from Mary represents the joy arising from the love of the conjunction of good and truth, thus the joy of heavenly conjugial love, which is in every particular of the Word. 

\Note 1:41 {filled with the Holy Ghost}
{\it Lord 51:5, TCR 158} The {\bf Holy Ghost} is the Divine proceeding, or the Divine Life emanating from the Word. This can also be expressed as the unified Divine Love and Wisdom proceeding from the Lord, expressed in someone leading a regenerate life. It is worth noting that the {\bf Holy Ghost} isn't mentioned in the Old Testament, but the Spirit of Holiness is mentioned three times. This is because the {\bf Holy Ghost} emanates from the Lord, and couldn't not proceed from Him until He had been born in this world.

\Note 1:46 {My soul}
{\it AE 750:5} The {\bf soul} has several meanings within the Word: in general, it means man and everything pertaining to him, more specifically it means a man's spiritual life, and more specific still it means the life of a man's understanding.

\Note 1:47 {spirit}
{\it Lord 49} The {\bf spirit} refers to the life of a regenerate person, or the spiritual life i.e. the life the Lord intends for us.

\Note 1:50 {them that fear him}
{\it AE 696:21, AR 527:2} The {\bf fear} of the Lord signifies having a sense of holiness and reverence and accordingly worshiping with holiness and reverence. It does not signify the same negative fearfullness we interpret in modern language, it signifies the opposite. "To fear" signifies to love because everyone who loves also fears to do evil to him whom he loves. Genuine love is not given without that fear.




















